\section{Introduction}
\label{sec:intro}

\IEEEPARstart{C}{onsider} an armed robbery registered at a town police
station. Three men rob a jewellery shop at closing time, an employee is
injured, and one robber leaves blood on a display counter. Within two weeks
the case file contains the first information report, the complainant and
witness statements, a medico-legal certificate, the investigating officer's
case diary, seizure memos, CCTV footage, forensic exhibits, and the call
records of a stolen phone. During investigation, one of the accused is found
to be seventeen years old, so his identity and records fall under statutory
juvenile protection. From this point, requests for the file arrive from
several directions, and each one needs a different answer.

The sub-inspector assigned as investigating officer needs the full file and
should receive it. A constable from the same station, with no connection to
the case, asks for the case diary and should be refused, although his station
and his role are the same as his colleague's. A sub-inspector from a
neighbouring district, investigating a similar gang robbery, requests the
CCTV stills to compare the method of operation; his purpose is legitimate,
but the request crosses a district boundary, so it should be neither silently
granted nor silently refused. It should go to a supervisory officer, and the
approval itself should be recorded. A forensic scientist analysing the blood
sample needs the exhibit-chain records and should get them, but when the same
scientist requests the victim's name and address a few minutes later, the
answer must be no, because nothing in a forensic examination requires it. The
superintendent of police, despite outranking everyone above, must be refused
the juvenile's identity, because a statutory protection class is stronger
than rank. Finally, six months later, an inquiry authority who was present
for none of these events must be able to verify what was released, to whom,
on what grounds, and who approved it.

This scenario is the subject of the paper, and it shows three things that any
access mechanism for such workflows must provide. First, the decision cannot
be derived from the job role alone: the constable and the investigating
officer hold comparable roles and receive opposite answers, and the forensic
scientist receives opposite answers to two of his own requests made minutes
apart. The decision depends on attributes of the requester, the record, the
requested action, the stated purpose, and the operating context. Second,
allow and deny are not enough: some requests are legitimate but cross a
boundary that requires a human authority, so \emph{escalate} must be a
first-class outcome and the resulting approval must become part of the
evidence. Third, decision evidence must outlive the decision: whatever the
inquiry authority needs in month six must have been recorded in month one,
because it cannot be reconstructed afterwards. The risk of weak governance
here is not hypothetical. In 2023, a data breach at the Police Service of
Northern Ireland exposed the names, ranks, and locations of serving officers,
and the UK Information Commissioner's Office later imposed a fine over the
incident~\cite{guardian_psni_breach_2023,ico_psni_fine_2024}.

The infrastructure to move such records between agencies already exists. In
India, the Crime and Criminal Tracking Network and Systems (CCTNS) connects
police stations nationwide, and the Inter-Operable Criminal Justice System
(ICJS) links the police pillar with courts, prisons, forensic laboratories,
and prosecution~\cite{pib_cctns_2026,mha_icjs_2026}. This paper does not
propose a replacement for these systems. It addresses what they leave open:
once a record is technically reachable across an agency boundary, the
decision to release it must be made from context, recorded with its grounds,
and protected against later manipulation.

Existing work covers parts of this requirement, and each part stops short at
a known point. Role-based access control assigns permissions to job
titles~\cite{sandhu_rbac_1996}, but the scenario above defeats it in its
first two requests. Attribute-based access control decides from subject,
object, action, and environment attributes~\cite{nist_abac_2014}, and recent
work extends it toward accountability~\cite{hu_accountable_private_ac_2024}
and privacy preservation~\cite{kerl_privacy_abac_2025}; policy engines of
this kind, however, are built to produce decisions, not to preserve
reviewable evidence of how each decision was reached. Blockchain-based audit
preserves evidence in a tamper-evident, multi-party
form~\cite{nist_blockchain_access_2022,androulaki_fabric_2018} and has been
applied to digital evidence
protection~\cite{banerjee_dashdikpala_2026,anbarasi_ai_criminal_justice_2026},
but it certifies only that a stored entry has not been altered after it was
written. Explainable AI produces reasons for decisions, and the argument that
high-stakes systems should prefer reviewable decisions is well
established~\cite{rudin_2019,zocholl_xai_law_enforcement_2025}; in practice,
however, explanations are rendered for the requester and discarded, rather
than stored as evidence bound to the decision they justify.

There is also a failure mode that none of the audit mechanisms above can
detect, and the scenario makes it concrete. Suppose the denial issued to the
forensic scientist for the victim's identity is later rewritten into an
approval by an attacker who controls the component that signs audit entries.
The rewritten entry is signed with the genuine key, so its signature
verifies, the chain remains intact, and every agency's copy of the ledger
agrees. The inquiry authority in month six runs an integrity check and finds
nothing wrong, yet the ledger now preserves a decision that policy never
made. Integrity and correctness are separate properties: a signature answers
whether the log was changed after writing, not whether what was written was
true.

This paper proposes SEBA-XAI, a Secure, Explainable, Blockchain-Audited
access-governance overlay that addresses the scenario end to end. For every
request, a policy service evaluates the attributes above and returns
\emph{allow}, \emph{deny}, or \emph{escalate}. At decision time, the system
records the grounds of the decision: the attributes that were decisive and
the policy version applied. The decision and its grounds are committed to a
permissioned blockchain audit trail shared by the participating agencies, so
that no single agency can silently rewrite history, while raw records never
leave the agency holding them. Because the re-signing attack defeats even
this audit trail from the inside, the paper additionally evaluates detection
methods that examine the recorded decisions rather than their signatures.

\subsection{Practical Challenges}

The paper is organised around four practical challenges taken directly from
the scenario. Each challenge is paired with one contribution in the same
order.

\begin{itemize}
    \item \textbf{Contextual access governance.} Access cannot be decided
    from static roles: the constable and the investigating officer share a
    role but not an assignment, and the forensic scientist's two requests
    differ only in the record and purpose. The decision must consider the
    requester, record, action, purpose, approval state, and operating
    context.

    \item \textbf{Tamper-evident audit across agencies.} The inquiry
    authority can trust the decision evidence only if no agency could have
    silently changed it, and raw sensitive records must not be copied to a
    shared ledger to achieve this.

    \item \textbf{Explainable access decisions.} Every allow, deny, or
    escalate outcome in the scenario must carry its grounds, because these
    decisions affect investigation timelines, privacy protection,
    inter-agency accountability, and later legal or administrative review.

    \item \textbf{Valid-looking corrupted logs.} An attacker who controls a
    signing component can rewrite a decision and re-sign it, so the audit
    trail remains internally valid while recording a decision that was never
    made.
\end{itemize}

\subsection{Research Contributions}

To address these challenges, the paper makes four contributions, ordered to
match the challenges above.

\begin{itemize}
    \item \textbf{Contribution 1:} A layered access-control service that
    combines role-based, attribute-based, and policy-based rules and returns
    \emph{allow}, \emph{deny}, or \emph{escalate} for inter-agency record
    requests, so that cases such as the constable's request and the forensic
    scientist's second request are decided from context rather than role.

    \item \textbf{Contribution 2:} A permissioned blockchain audit layer in
    which each decision event records the request identifier, decision hash,
    explanation hash, policy version, approval reference where one exists,
    and record commitment, while raw records remain off-chain with the
    holding agency.

    \item \textbf{Contribution 3:} An explainable decision service that
    stores the decisive attributes, reason code, policy version, and
    counterfactual information as a structured artifact hash-linked to the
    audit event, so that a reviewer can verify the grounds of a decision
    without opening the underlying records.

    \item \textbf{Contribution 4:} A reproducible adversarial benchmark
    covering ordinary tampering attacks and the validly re-signed
    compromised-signer attack. Across five seeds of a synthetic workload,
    every integrity-based defence tested detects ordinary tampering but
    obtains 0.00 detection for the re-signed attack, while two methods that
    examine the decisions themselves, a statistical drift detector over the
    decision log and a trusted attribute oracle, obtain 1.00 detection under
    stronger visibility assumptions.
\end{itemize}

The scope of the claims is deliberately limited. Real police access logs and
records are not publicly available, so the evaluation uses a reproducible
synthetic workload, and the blockchain layer is a simulated permissioned
ledger rather than a live deployment. The paper does not claim deployment
inside any operational system, legal compliance, validation on real police
records, or production security. It also deliberately avoids crime
prediction: aggregate public crime statistics are not individual-level
access-control data~\cite{ncrb_2023}, predictive policing can create
feedback loops when observed crime data is shaped by earlier policing
decisions~\cite{ensign_feedback_2018}, and complex criminal-justice
prediction systems do not reliably outperform simple
baselines~\cite{dressel_fairness_limits_2018}. The subject of this paper is
the governance of access, not the prediction of crime.

The remainder of the paper is organised as follows.
Section~\ref{sec:related} reviews related work.
Section~\ref{sec:methodology} presents the SEBA-XAI architecture, threat
model, and benchmark design. Section~\ref{sec:results} reports the results,
including the compromised-signer asymmetry.
Section~\ref{sec:limitations} states the limitations, and
Section~\ref{sec:conclusion} concludes.
